%% =====================================================================
%% FeatureLiftBench -- FSE submission draft
%% Based on the user-provided ACM acmart template.
%% Double-blind review mode:
%%   \documentclass[acmsmall,anonymous,review,screen]{acmart}
%% Camera-ready mode:
%%   \documentclass[acmsmall]{acmart}
%% =====================================================================

\documentclass[acmsmall,anonymous,review,screen]{acmart}

%TC:macro \cite [option:text,text]
%TC:macro \citep [option:text,text]
%TC:macro \citet [option:text,text]
%TC:envir table 0 1
%TC:envir table* 0 1
%TC:envir tabular [ignore] word
%TC:envir displaymath 0 word
%TC:envir math 0 word
%TC:envir comment 0 0

\usepackage{booktabs}
\usepackage{listings}
\usepackage{tabularx}
\usepackage{longtable}
\usepackage{xspace}

\lstset{
  basicstyle=\ttfamily\small,
  keywordstyle=\bfseries,
  commentstyle=\itshape\color{gray},
  breaklines=true,
  frame=single,
}

\newcommand{\flb}{\textsc{FeatureLiftBench}\xspace}
\newcommand{\rres}{\textsc{RRES}\xspace}
\newcommand{\mainprotocol}{\textsc{Official Main}\xspace}
\newcommand{\draftnote}[1]{\textcolor{red}{[Draft note: #1]}}

%% Submission-stage placeholders. Replace from the ACM rights form for camera ready.
\setcopyright{acmlicensed}
\copyrightyear{2026}
\acmYear{2026}
\acmDOI{XXXXXXX.XXXXXXX}
\acmConference[FSE '26]{Proceedings of the ACM International Conference
  on the Foundations of Software Engineering}{September 14--18, 2026}{Seoul, South Korea}
\acmISBN{978-1-4503-XXXX-X/26/09}
%% \acmSubmissionID{TBD}

\begin{document}

\title{FeatureLiftBench: Evaluating Repository-Level Feature Extraction by Code Agents}

%% The anonymous class hides this block during review. Replace for camera ready.
\author{Anonymous Author(s)}
\affiliation{%
  \institution{Anonymous Institution}
  \city{Anonymous City}
  \country{Anonymous Country}
}
\renewcommand{\shortauthors}{Anonymous Author(s)}

\begin{abstract}
Extracting a reusable feature from an existing repository differs from fixing
a reported issue or generating a new implementation from scratch. An agent
must locate relevant implementation evidence, recover its transitive behavioral
and resource closure, adapt it into a standalone package, and preserve
observable behavior without retaining a runtime dependency on the source
repository. Existing code-agent benchmarks do not isolate this combination of
repository understanding, behavioral preservation, modularization, and
compactness. We introduce \flb, a benchmark for repository-level,
behavior-preserving feature extraction by code agents. Its main Python suite
contains 200 tasks drawn from 176 pinned open-source repositories. Under the
Full-Repository / No-Hint protocol, an agent receives the complete repository
and a complete public behavioral contract, but no source-location hints,
benchmark tests, Hidden tests, or reference implementation. Deterministic
Dockerized checks evaluate buildability, Public behavior, deeper Hidden
behavior, and independence from the source repository. We separately measure
the compactness of functionally passing packages relative to evaluator-side
feasible references. Under a frozen OpenHands execution protocol, current
agents achieve \draftnote{insert the final eligible Python-200-prime
cross-model range}. Artifact and trajectory analyses identify recurring
contract-closure failures, in which plausible packages omit required exports,
exception semantics, state transitions, resources, or preservation details.
Controlled information experiments further show that executable Public
feedback can repair Public failures without generally eliminating Hidden
failures, while evaluated self-testing, repair, checkpointing, and budget
scaffolds do not stably outperform the strongest legal Main protocol. \flb
provides a reproducible basis for measuring how code agents turn repository
evidence into independent, behavior-complete artifacts.
\end{abstract}

\begin{CCSXML}
<ccs2012>
 <concept>
  <concept_id>10011007.10011074.10011099.10011102</concept_id>
  <concept_desc>Software and its engineering~Software testing and debugging</concept_desc>
  <concept_significance>500</concept_significance>
 </concept>
 <concept>
  <concept_id>10011007.10011074.10011111.10011113</concept_id>
  <concept_desc>Software and its engineering~Software verification and validation</concept_desc>
  <concept_significance>300</concept_significance>
 </concept>
 <concept>
  <concept_id>10011007.10011074.10011099.10011693</concept_id>
  <concept_desc>Software and its engineering~Software maintenance tools</concept_desc>
  <concept_significance>300</concept_significance>
 </concept>
</ccs2012>
\end{CCSXML}

\ccsdesc[500]{Software and its engineering~Software testing and debugging}
\ccsdesc[300]{Software and its engineering~Software verification and validation}
\ccsdesc[300]{Software and its engineering~Software maintenance tools}

\keywords{code agents, repository-level evaluation, feature extraction,
behavior preservation, software modularization, benchmark}

\maketitle

\section{Introduction}
\label{sec:introduction}

Software maintenance frequently requires engineers to extract an existing
capability from a large system: a parser must become a small library, a
configuration resolver must be reused in another service, or a plugin registry
must be separated from the framework that originally hosted it. This operation
is common in modernization, dependency reduction, service decomposition, and
reuse of legacy code. It is also deceptively difficult. The target behavior may
be distributed across helper functions, stateful registries, resources,
configuration defaults, error classes, and third-party dependencies. Copying
the most obvious implementation region may produce a package that works on a
happy path while silently changing edge behavior or retaining hidden
dependencies on the original repository.

Repository-level code agents appear well suited to this task. They can search
large workspaces, inspect call sites and tests, edit multiple files, execute
tools, and package new artifacts. However, current evaluation paradigms do not
isolate whether an agent can perform behavior-preserving feature extraction.
Issue-resolution benchmarks ask an agent to repair a repository until tests
pass~\cite{jimenez2024swebench}. Greenfield code-generation benchmarks ask for
a new implementation from a specification~\cite{chen2021codex}.
Repository-understanding and localization studies test whether the relevant
code can be found. None of these settings simultaneously requires an agent to
recover an existing feature from an entangled repository, preserve a declared
behavioral contract, remove dependence on the original package, and avoid
solving the task through broad repository copying.

Feature extraction therefore combines at least three obligations:
\begin{enumerate}
  \item \textbf{Locate:} identify implementation evidence relevant to the
        requested capability;
  \item \textbf{Close:} recover the APIs, helpers, state, resources,
        dependencies, and edge semantics needed for behavioral completeness;
  \item \textbf{Isolate:} package the recovered feature so that it executes
        independently of the source repository.
\end{enumerate}
We use this decomposition as an explanatory model rather than a causal
factorization. A final failure does not by itself reveal which obligation
failed; semantic labels require trajectory and artifact evidence.

We introduce \flb, a benchmark that evaluates these obligations under a
controlled Full-Repository / No-Hint protocol. Each task provides a complete
repository pinned to an immutable source revision and a complete public
contract describing the required package interface and observable behavior.
The agent may inspect the entire repository, but it is not given source-location
hints, benchmark tests, Hidden tests, or a reference solution. It must produce
a standalone \texttt{featurelifted} package. A submission passes only when it
builds, satisfies both primary and deeper contract tests, and remains functional
after access to the original repository is removed.

\flb separates two quality dimensions. \emph{Functional Pass} measures whether
the submission satisfies build, Public, Hidden, and isolation gates.
\emph{Reference-Relative Extraction Size} (\rres) measures the normalized
footprint of a functionally passing package relative to an evaluator-side
feasible reference. Correctness and compactness are not combined into one
score: a broad vendoring solution may preserve behavior but fail to demonstrate
a compact extraction, while a small package may simply be incomplete.

The main Python-200$'$ suite combines a frozen 150-task baseline with a
separately calibrated Hard-50 split, for 200 tasks from 176 repositories. The
Hard-50 split increases coverage of plugin registries, lifecycle and session
behavior, multi-source configuration, validation boundaries, deep parsing, and
direct copy-trap tasks. An earlier External-50 expansion is retained only as an
easy, copy-heavy side split: strong-agent pass rates of 90--94\% and
pass-conditioned footprints close to whole-repository copying showed that it
was unsuitable for the paper's main difficulty claim.

Our empirical study asks four groups of questions:
\begin{itemize}
  \item \textbf{Capability:} How often do current code agents produce
        independent, behavior-complete feature packages?
  \item \textbf{Quality:} When agents pass, how compact are their extractions?
  \item \textbf{Failure mechanisms:} At which gate do agents fail, and which
        contract obligations remain open?
  \item \textbf{Information and resources:} How do executable feedback,
        process scaffolds, and trajectory cost change outcomes?
\end{itemize}

Our analysis points to a recurring gap between plausible repository work and
verifiable feature closure. Strong agents frequently produce buildable packages
and pass primary behavior checks, yet fail on required exports, boundary
conditions, exception semantics, stateful behavior, resource closure, or
preservation details. We refer to this recurring symptom as
\emph{contract-closure failure}. It is not inferred from a Hidden failure alone:
attribution requires evidence that the relevant requirement was public and that
the trajectory or artifact omitted or changed it. Controlled information
experiments further suggest that executable Public feedback and Hidden
behavioral completeness are distinct layers.

This paper makes three main contributions:
\begin{itemize}
  \item \textbf{Benchmark asset.} We define repository-level,
        behavior-preserving feature extraction and construct a suite of 200
        tasks from 176 pinned Python repositories, with public contracts,
        source registries, deterministic evaluators, evaluator-side feasible
        references, isolation checks, and task-level provenance. A v2
        contract-completeness label on freeze~v2 is 200/0/0 after repairing
        32 confirmed surface or entry-point defects recorded on the
        predecessor freeze (then 168/32/0), without changing suite membership.
  \item \textbf{Evaluation protocol.} We introduce a Full-Repository / No-Hint
        protocol that fixes repository, contract, runtime, budget, and evaluator
        while varying model backends. We report deterministic Functional Pass
        separately from pass-conditioned compactness and process cost.
  \item \textbf{Diagnostic analysis.} We analyze failure stages, semantic
        contract-closure symptoms, task factors, information boundaries, and
        trajectory cost. We also report negative results from legal
        self-testing, repair, checkpoint, context, and budget-control
        interventions rather than presenting them as successful methods.
\end{itemize}

\draftnote{After Gate A and Gate C, add one sentence with the eligible
Python-200-prime capability range and the dominant failure-stage result.}


\section{Related Work}
\label{sec:related-work}

\subsection{Executable Software-Engineering Benchmarks}

Executable software-engineering benchmarks have established the importance of
real repositories, tool use, sandboxed execution, and test-based verification.
SWE-bench asks language models to resolve real GitHub issues in existing
repositories~\cite{jimenez2024swebench}, while Terminal-Bench evaluates agents
on realistic command-line tasks~\cite{merrill2026terminalbench}. Their central
unit is typically an issue, a failing test, or a requested modification to the
original workspace. \flb instead evaluates extraction into a new independent
artifact. The source repository is evidence, not the final execution
environment, and passing requires behavior preservation after the source
package is unavailable.

Agent-computer interfaces and runtime design can also change executable-agent
performance~\cite{yang2024sweagent,yao2026harnessbench}. We therefore fix
OpenHands as the \mainprotocol runtime for the model leaderboard. Runtime
comparisons, when available, are reported separately rather than mixed into
model scores.

\subsection{Code Generation, Repository Understanding, and Localization}

Function-level code generation measures whether a model can synthesize a
program from a natural-language or executable specification, as in HumanEval
and related synthesis
suites~\cite{chen2021codex,austin2021mbpp,lai2023ds1000,zhuo2024bigcodebench,jain2024livecodebench}.
Repository-level completion and retrieval study whether a model can use
in-repo context to finish or locate code~\cite{zhang2023repocoder,liu2024repobench}.
Feature-location research asks whether a developer or tool can identify files,
entities, or regions relevant to a concern~\cite{dit2013featurelocation}.
Feature lifting requires both inference and artifact construction, but success
is not reducible to either. Locating the correct implementation is insufficient
if transitive behavior or resources are omitted; generating a behaviorally
similar implementation is insufficient if the public contract or independent
package boundary is not preserved.

\subsection{Program Slicing, Modularization, and Library Extraction}

Feature extraction is related to program slicing~\cite{weiser1984slicing},
multi-dimensional separation of concerns~\cite{tarr1999degrees}, and the
refactoring literature that surveys how developers extract and modularize
behavior~\cite{mens2004refactoring}. Classical techniques typically assume a
program representation, slicing criterion, or human-specified boundary. \flb
evaluates an autonomous agent from a natural-language public contract and a
complete repository, and grades the final independent artifact by observable
behavior. The benchmark does not claim to compute a minimal semantic slice: its
frozen references are feasible comparison points, and \rres is a compactness
proxy rather than a proof of optimality.

\subsection{Information Boundaries in Agent Evaluation}

Agent performance depends on the model, execution runtime, tools, context
policy, and feedback interface~\cite{yang2024sweagent,wang2024openhands,yao2026harnessbench}.
\flb separates these factors in two ways. The main leaderboard fixes OpenHands
as the \mainprotocol runtime and varies the model backend. Information
ablations instead change what task evidence is visible while preserving model,
runtime, evaluator, and budget. This design distinguishes model comparison from
questions such as whether executable Public tests or source-location hints
alter performance. Runtime replications, when completed, are reported
separately rather than mixed into model scores.


\section{The FeatureLiftBench Benchmark}
\label{sec:benchmark}

\subsection{Repository-Level Feature-Lifting Setting}

A task consists of a pinned source repository $R$, a public contract $S$, an
initial workspace $W$, and a private evaluator $J$. Given $(R,S,W)$, an agent
produces a package $P$ under a fixed execution protocol:
\begin{equation}
  P = \operatorname{Run}(M, H, R, S, W),
  \label{eq:run}
\end{equation}
where $M$ is the model backend and $H$ is the fixed \mainprotocol runtime. The
evaluator observes the collected package and assigns gate-level outcomes:
\begin{equation}
  \operatorname{Functional}(P)
  = B(P) \land P_{\mathrm{pub}}(P)
    \land H_{\mathrm{hid}}(P) \land I(P).
  \label{eq:functional}
\end{equation}
Here, $B$ checks installation and imports, $P_{\mathrm{pub}}$ checks primary
behavior from the public contract, $H_{\mathrm{hid}}$ checks deeper
combinations and boundary behavior from that same contract, and $I$ checks
independence from the source repository. A valid Hidden test must map to a
published contract requirement.

The agent receives the complete source tree and contract but no source-location
hint. It does not see Public tests, Hidden tests, evaluator code, or the frozen
reference. Extraction, adaptation, and behaviorally equivalent reimplementation
are allowed. Runtime imports from the source project, forbidden paths, and
undeclared dependency channels are disallowed.

\begin{figure}[t]
  \centering
  \fbox{%
    \parbox{0.92\linewidth}{\centering
      Pinned repository + public contract\\[2pt]
      $\downarrow$\\[2pt]
      Full-Repository / No-Hint agent execution\\[2pt]
      $\downarrow$\\[2pt]
      Standalone \texttt{featurelifted} package + trajectory\\[2pt]
      $\downarrow$\\[2pt]
      Build $\rightarrow$ Public $\rightarrow$ Hidden $\rightarrow$ Isolation\\[2pt]
      $\downarrow$\\[2pt]
      Functional Pass + pass-conditioned compactness}}
  \caption{The \flb evaluation pipeline. Protected tests, evaluator code, and
  evaluator-side references are never exposed during agent execution.}
  \Description{A pipeline from a pinned repository and public contract through
  agent execution and a standalone package to build, Public, Hidden, and
  isolation evaluation, followed by functional and compactness reporting.}
  \label{fig:pipeline}
\end{figure}

\subsection{Task Suite Design}

The paper suite, Python-200$'$, contains 200 tasks from 176 repositories
(Table~\ref{tab:suite}). It emphasizes Python libraries, developer tools, and
framework components. One source repository may support multiple distinct
tasks in the baseline split. Every task pins its upstream revision and records
source identity in a canonical registry.

\begin{table}[t]
  \caption{Composition of the Python-200$'$ paper suite.}
  \label{tab:suite}
  \centering
  \begin{tabular}{lrrl}
    \toprule
    \textbf{Split} & \textbf{Tasks} & \textbf{Repositories} & \textbf{Role} \\
    \midrule
    Frozen Python-150 & 150 & 127 & Baseline set \\
    Hard-50 & 50 & 50 & Calibrated expansion \\
    \midrule
    Python-200$'$ & 200 & 176 & Main paper suite \\
    \bottomrule
  \end{tabular}
\end{table}

Tasks are organized along three descriptive axes. \emph{Feature family}
records requested behavior, such as parsing, serialization, configuration,
validation, resource handling, registry behavior, caching, protocols,
algorithms, or workflows. \emph{Primary entanglement} records the dominant
mechanism that makes independent extraction difficult: data-model,
parser-state, framework, configuration/environment, resource, or third-party
dependency entanglement. \emph{Lift type} describes the relationship between
the requested artifact and source implementation: Direct, Adapted, or
Composite. These labels support sampling and analysis; they are not shown to
the agent and do not affect Functional Pass.

The combined suite contains 68 Direct, 100 Adapted, and 32 Composite tasks.
The Hard-50 expansion increases coverage of registry/plugin dispatch, workflow
and session orchestration, configuration discovery, validation boundaries,
deep parsing, and direct tooling copy traps. It was calibrated with a strong
model at 29/50 functional passes, within a predeclared 40--65\% target band.
This is benchmark-design evidence, not a substitute for the final main-table
evaluation.

An earlier External-50 expansion is retained as a side split rather than part
of Python-200$'$. Strong-model pass rates of 90--94\% and pass-conditioned
footprints close to one reference-relative unit revealed an easy, copy-heavy
construction. This negative design result motivated Hard-50 and illustrates why
functional outcome and extraction compactness must be measured separately.

\subsection{Task Construction and Validation}

Each candidate task is retained only when it satisfies four requirements:
\begin{enumerate}
  \item \textbf{Realism:} the requested feature corresponds to a plausible
        reuse or modularization need in a real repository;
  \item \textbf{Solvability:} a feasible independent implementation can be
        produced from the pinned repository and public contract;
  \item \textbf{Oracle-checkability:} success can be verified with
        deterministic build, behavior, and isolation checks;
  \item \textbf{Integrity:} the agent cannot obtain credit by reading
        protected tests, retaining a forbidden source dependency, or bypassing
        the intended artifact boundary.
\end{enumerate}

Task construction records a source revision and digest, public contract,
required API, locked dependencies, Public/Hidden contract mapping,
deterministic evaluator, and a feasible reference. Evaluator-side references
are never exposed to agents. Automated and manual gates check source
materialization, package imports, reference behavior, test isolation,
forbidden dependencies, resource access, and freeze identity.

The benchmark separates task packages from mutable working-tree state. Paper
results must be evaluated against the frozen materialization identified by the
active suite manifest. A source mismatch, unavailable locked dependency,
evaluator defect, or context-policy violation is an eligibility or
infrastructure outcome rather than a model failure.

\subsection{Run Protocol and Evidence Collection}

Each run follows a setup--execution--evaluation pipeline
(Figure~\ref{fig:pipeline}). Setup materializes the pinned task workspace and
records task, source, model, runtime, prompt, budget, and container identity.
During execution, the agent may inspect the full repository, use shell and
editing tools, run legal self-tests, and modify its submission workspace. The
runner collects the final submission tree, model usage, actions, errors, and
workspace snapshots. Evaluation installs the submission in a Dockerized
offline environment and applies build, Public, Hidden, and isolation checks.

Each task attempt produces four evidence layers: (1) the final artifact,
(2) the agent trajectory and process status, (3) token, step, time, context,
and resource statistics, and (4) the task-level deterministic evaluator result.
The evaluator result is authoritative for Functional Pass. Agent completion or
runner status is reported separately because an agent may hit a step limit after
leaving a passing package, or may report completion while its artifact fails.

\subsection{Metrics}

\paragraph{Functional Pass@1.}
A task receives one functional point only if all gates in
Equation~\ref{eq:functional} pass. Functional failures receive one exclusive
first-failure stage:
\begin{equation}
  \text{missing} \rightarrow \text{build} \rightarrow \text{public}
  \rightarrow \text{hidden} \rightarrow \text{isolation}.
\end{equation}
Infrastructure and evidence-availability flags are reported separately from
semantic model outcomes.

\paragraph{Compactness.}
For a functionally passing submission $P$ with evaluator-side feasible
reference $P_{\mathrm{ref}}$, we report
\begin{equation}
  \operatorname{RRES}(P)
  = \frac{\operatorname{normalized\_size}(P)}
         {\operatorname{normalized\_size}(P_{\mathrm{ref}})}.
  \label{eq:rres}
\end{equation}
\rres is accompanied by file, copied-code, and dependency-footprint
diagnostics where evidence is available. It is reported only on functional
passes. Cross-method compactness comparisons use the same-task subset on which
both methods pass.

\paragraph{Process and cost.}
Tokens, steps, latency, completion status, and resource failures are diagnostic
metrics. On passing trajectories with replayable package snapshots, $T^*$ is
the earliest cumulative token point at which a unique package tree passes the
evaluator. $T^*$ is undefined for failing trajectories and is not itself a
deployable stopping rule.

\section{Experiments}
\label{sec:experiments}

\subsection{Experimental Setup}

The headline study evaluates five model backends on the frozen Python-150
split of Python-200$'$ under freeze~v2
(\texttt{6c20ff0307762503a73cbb9ff32e9992c6446e4b17483a68373027be58cbf419}).
Table~\ref{tab:factors} fixes the execution boundary. The Official Hard-50
split is part of the released 200-task suite and is reported in
Appendix~\ref{app:hard50}; it is not the headline leaderboard.

\begin{table*}[t]
  \caption{Controlled and varying factors in the \mainprotocol evaluation.}
  \label{tab:factors}
  \centering
  \begin{tabularx}{\textwidth}{@{}p{0.34\textwidth}X@{}}
    \toprule
    \textbf{Factor} & \textbf{Official Main treatment} \\
    \midrule
    Task contract and repository snapshot & Fixed per task (freeze~v2) \\
    Initial workspace and source registry & Fixed per task \\
    Agent runtime & OpenHands \mainprotocol \\
    Prompt, action budget, and context envelope
      & Fixed (120 steps, 128k envelope, no source-location hint) \\
    Evaluator and Docker images
      & \texttt{featureliftbench-agent/eval:python200-prime-212930ea} \\
    Model backend & Varied in the main matrix \\
    Public-test visibility & Withheld in Main \\
    Source-location hint & Withheld in Main \\
    Evaluator, reference, and Hidden tests & Never exposed to the agent \\
    \bottomrule
  \end{tabularx}
\end{table*}

Each cell is one attempt per task. Functional Pass is
$B \land P_{\mathrm{pub}} \land H_{\mathrm{hid}} \land I$
(Equation~\ref{eq:functional}). An empty submission is a functional failure.
Agent or runner \texttt{status} is not the score: a package may pass the
evaluator after a non-zero process code, and a ``passed'' runner status is
not required. We report Wilson 95\% intervals for rates and exact McNemar
tests on discordant tasks. One attempt does not estimate run-to-run variance.

The five backends are DeepSeek V4 Pro, DeepSeek V4 Flash, gpt-5.6-luna
served via OpenLux, Qwen3.6-35B-A3B-FP8, and GPT-OSS 120B. Pro was run on
the Python-150 suite directory; the other four were run on the 200-task
Official Main directory and then restricted to the same 150 task IDs. Runs
were collected on 2026-09-04/05. An unfinished GLM-5.3-Flash suite is
excluded. Configuration digests are in Appendix~\ref{app:configuration}.

\subsection{Overall Capability}
\label{sec:capability}

Table~\ref{tab:main} is the paper main table: Functional Pass@1 on frozen
Python-150. Core-100 and hard3 are the in-suite construction split of those
150 tasks, not the official Hard-50 expansion.

\begin{table*}[t]
  \caption{Functional Pass@1 on freeze~v2 Python-150 (Official Main,
  OpenHands). Empty submissions count as failures. Core-100 / hard3 is the
  in-suite construction split.}
  \label{tab:main}
  \centering
  \begin{tabular}{lrrrrrr}
\toprule
\textbf{Model} & \textbf{Pass} & \textbf{Rate} & \textbf{Wilson 95\%}
  & \textbf{Core-100} & \textbf{hard3} & \textbf{Empty} \\
\midrule
DeepSeek V4 Pro
  & 115/150 & 76.7\% & [69.3, 82.7]
  & 94/100 & 21/50 & 0 \\
DeepSeek V4 Flash
  & 108/150 & 72.0\% & [64.3, 78.6]
  & 90/100 & 18/50 & 0 \\
gpt-5.6-luna (OpenLux)
  & 102/150 & 68.0\% & [60.2, 74.9]
  & 82/100 & 20/50 & 6 \\
Qwen3.6-35B
  & 63/150 & 42.0\% & [34.4, 50.0]
  & 55/100 & 8/50 & 25 \\
GPT-OSS 120B
  & 36/150 & 24.0\% & [17.9, 31.4]
  & 25/100 & 11/50 & 2 \\
\bottomrule
\end{tabular}

\end{table*}

Three facts follow from the table. First, the benchmark is not saturated:
the strongest configuration still fails 35/150 tasks, and 28 tasks are
unsolved by all five models (22 of them hard3). Second, backends separate
clearly. Pro exceeds Qwen by 34.7 points and OSS by 52.7 points (exact
McNemar $p\ll 0.001$; Table~\ref{tab:python150-v2-pairwise}). Pro versus Luna
is 18/5 discordant tasks ($p=0.011$). Pro versus Flash is 10/3 ($p=0.092$)
with overlapping Wilson intervals, so we do not claim that Pro is
significantly stronger than Flash. Flash versus Luna is not significant
($p=0.33$). Third, Qwen's 25 empty submissions remain functional failures;
we do not rewrite that row as 63/125.

\paragraph{Finding 1.}
Current coding agents exhibit substantial but incomplete feature-lifting
capability on frozen Python-150, with large performance differences across
model backends.

\subsection{Failure Stages}
\label{sec:stages}

Table~\ref{tab:funnel} reports mutually exclusive first outcomes. Pro and
Flash never fail to deliver a package and never fail Build. Across the
matrix, Isolation-first failure occurs twice. The mass of residual failure
is Public then Hidden. OSS fails Public 70 times; Qwen's 25 missing
submissions are a process event and are separated in
Section~\ref{sec:process}.

\begin{table*}[t]
  \caption{Mutually exclusive first outcomes on freeze~v2 Python-150.}
  \label{tab:funnel}
  \centering
  \begin{tabular}{lrrrrrr}
\toprule
\textbf{Model} & \textbf{Pass} & \textbf{Missing} & \textbf{Build}
  & \textbf{Public} & \textbf{Hidden} & \textbf{Isolation} \\
\midrule
DeepSeek V4 Pro & 115 & 0 & 0 & 25 & 10 & 0 \\
DeepSeek V4 Flash & 108 & 0 & 0 & 26 & 15 & 1 \\
gpt-5.6-luna (OpenLux) & 102 & 6 & 3 & 27 & 12 & 0 \\
Qwen3.6-35B & 63 & 25 & 6 & 35 & 21 & 0 \\
GPT-OSS 120B & 36 & 2 & 18 & 70 & 23 & 1 \\
\bottomrule
\end{tabular}

\end{table*}

\paragraph{Finding 2.}
Functional failures concentrate primarily at the behavioral gates: producing
a buildable package is substantially easier than recovering complete required
behavior.

\subsection{Process versus Capability}
\label{sec:process}

Functional Pass does not change with process status. Some failures nevertheless
occur before a gradable package exists. Table~\ref{tab:process} counts empty
submissions separately from later artifact failures. Qwen's 25 empty rows are
OpenHands \texttt{security\_risk} tool-validation errors, typically within
tens of seconds; they do not demonstrate Hidden semantic failure. Luna has
six empty rows (three tool-validation errors, two invalid encrypted
content events, and one file-less workspace). Flash records 82 functional
passes whose runner status is not \texttt{passed}; those rows remain passes.

\begin{table}[t]
  \caption{Empty submissions on Python-150. Empty rows are functional
  failures and are excluded from Section~\ref{sec:mechanism}.}
  \label{tab:process}
  \centering
  \begin{tabular}{lrr}
\toprule
\textbf{Model} & \textbf{Empty} & \textbf{TVE/encrypted empty} \\
\midrule
DeepSeek V4 Pro & 0 & 0 \\
DeepSeek V4 Flash & 0 & 0 \\
gpt-5.6-luna (OpenLux) & 6 & 5 \\
Qwen3.6-35B & 25 & 25 \\
GPT-OSS 120B & 2 & 2 \\
\bottomrule
\end{tabular}

\end{table}

Process failures affect some model--harness configurations, but they are
analytically distinct from artifact-level feature-lifting failures. Empty
submissions are excluded from Section~\ref{sec:mechanism}.

\section{Analysis}
\label{sec:analysis}

\subsection{Failure Mechanism}
\label{sec:mechanism}

Mechanical stages locate the first failing gate. They do not name a cause.
We therefore close-read artifact-level failures: the public contract, the
first-failure evaluator log, and the submitted package, plus a trajectory
screen for whether the agent inspected \texttt{repo/}. Labels follow the
protocol in Table~\ref{tab:symptoms}. Hidden test names, inputs, and
assertions are not reported. Coding is an assistant first pass (L1), not
independent human dual review, and is not a gold inter-rater study.

The census is a close-read of Pro and Flash artifact-level failures
($n=63$; Pro 28, Flash 35). We do not merge other backends into this
denominator.

\begin{table*}[t]
  \caption{Semantic symptoms used for artifact and trajectory analysis.}
  \label{tab:symptoms}
  \centering
  \begin{tabularx}{\textwidth}{@{}p{0.23\textwidth}Xp{0.27\textwidth}@{}}
    \toprule
    \textbf{Symptom} & \textbf{Observable manifestation} & \textbf{Required evidence} \\
    \midrule
    Localization & No inspection of the relevant source region
      & Trajectory + package \\
    Contract/API completion & Required symbol, member, or branch is absent
      & Public contract + package \\
    Dependency / resource closure & Helper, registry, config, or resource omitted
      & Source evidence + failing case \\
    Behavioral drift & API present; return, exception, or edge semantics differ
      & Contract + evaluator + package \\
    Packaging / modularization & Logic exists but is not independently installable
      & Package tree + isolation evidence \\
    Copy-heavy pass & Gates pass through broad vendoring
      & Passing result + RRES/copy evidence \\
    \bottomrule
  \end{tabularx}
\end{table*}

\begin{table}[t]
  \caption{Primary semantic cause on Pro and Flash artifact failures
  (L1 close-read, $n=63$). Not independent human gold.}
  \label{tab:mechanism}
  \centering
  \begin{tabular}{lrrr}
\toprule
\textbf{Primary cause} & \textbf{Total} & \textbf{Pro} & \textbf{Flash} \\
\midrule
Behavior drift & 54 & 26 & 28 \\
Contract/API completion & 7 & 2 & 5 \\
Packaging / modularization & 2 & 0 & 2 \\
Dependency / resource closure & 0 & 0 & 0 \\
Localization & 0 & 0 & 0 \\
Unknown & 0 & 0 & 0 \\
\midrule
Coded artifact failures & 63 & 28 & 35 \\
\bottomrule
\end{tabular}

\end{table}

Closure classes (API completion plus behavior drift) account for 61/63 rows.
The modal label is behavior drift: the required surface is present, but
observable semantics are not. Localization is 0 after the trajectory screen:
every remaining failure issued a deep read of \texttt{repo/} (find, grep,
cat, or file view). Packaging appears twice on Flash (an isolation
\texttt{forbidden\_imports} hit, and a submission that still imports the
upstream package after the isolation gate). We do not treat localization as
solved suite-wide; the claim is that, on this Pro+Flash artifact-fail slice,
the dominant residue is incomplete contract recovery rather than missing
packages or uninspected source trees.

A stratified L1 sample of Luna, Qwen, and OSS artifact failures
($n=45$ draws, 32 coded rows; seed
\texttt{python150-postsample-v1}) is reported in
Table~\ref{tab:postsample} and is \emph{not} pooled into the 63.
Luna remains in the same direction (drift and missing exports). Qwen and OSS
additionally fail Build by omitting transitive dependencies or by importing
the original package. That is a weaker-backend packaging story, not a
five-model pie chart.

\begin{table}[t]
  \caption{Stratified L1 sample of Luna/Qwen/OSS artifact failures
  (32 coded rows). Not merged with Table~\ref{tab:mechanism}.}
  \label{tab:postsample}
  \centering
  \begin{tabular}{lrrrr}
\toprule
\textbf{Primary cause} & \textbf{Total} & \textbf{Luna} & \textbf{Qwen} & \textbf{OSS} \\
\midrule
Behavior drift & 14 & 4 & 7 & 3 \\
Contract/API completion & 9 & 4 & 0 & 5 \\
Dependency / resource closure & 4 & 0 & 3 & 1 \\
Packaging / modularization & 5 & 0 & 1 & 4 \\
Localization & 0 & 0 & 0 & 0 \\
\midrule
Coded sample rows & 32 & 8 & 11 & 13 \\
\bottomrule
\end{tabular}

\end{table}

\paragraph{Finding 3.}
On the Pro+Flash artifact-fail slice, failures are dominated by incomplete
recovery of the required behavioral contract (a contract-closure gap), rather
than by inability to emit an installable package or by failing to inspect the
source tree. A stratified sample of the other backends does not reverse that
direction on Luna, but adds packaging and missing-helper failures on the
weaker backends. The proportions are an L1 first pass, not gold.

\subsection{Task Difficulty}
\label{sec:difficulty}

Table~\ref{tab:construction} splits Python-150 by construction cohort.
For Pro, Core-100 is 94/100 and hard3 is 21/50. Flash is 90/100 versus
18/50. Of 28 tasks unsolved by every model, 22 are hard3. OSS barely moves
(25\% versus 22\%). Entanglement level is uniformly high on this split, so
it cannot be tested. Lift type is directional but weaker
(Table~\ref{tab:lift}): Pro Direct 51/56, Adapted 55/76, Composite 9/18
(the last cell is small). Qwen's Adapted cell includes 16 empty submissions,
which are process events rather than a lift-type effect.

\begin{table}[t]
  \caption{Functional Pass by in-suite construction split on Python-150.
  hard3 is not the official Hard-50 expansion.}
  \label{tab:construction}
  \centering
  \begin{tabular}{lrr}
\toprule
\textbf{Model} & \textbf{Core-100} & \textbf{hard3} \\
\midrule
DeepSeek V4 Pro & 94/100 (94.0\%) & 21/50 (42.0\%) \\
DeepSeek V4 Flash & 90/100 (90.0\%) & 18/50 (36.0\%) \\
gpt-5.6-luna (OpenLux) & 82/100 (82.0\%) & 20/50 (40.0\%) \\
Qwen3.6-35B & 55/100 (55.0\%) & 8/50 (16.0\%) \\
GPT-OSS 120B & 25/100 (25.0\%) & 11/50 (22.0\%) \\
\bottomrule
\end{tabular}

\end{table}

\begin{table}[t]
  \caption{Functional Pass by lift type on Python-150. Composite $n=18$.}
  \label{tab:lift}
  \centering
  \begin{tabular}{lrrr}
\toprule
\textbf{Model} & \textbf{Direct} & \textbf{Adapted} & \textbf{Composite} \\
\midrule
DeepSeek V4 Pro & 51/56 (91.1\%) & 55/76 (72.4\%) & 9/18 (50.0\%) \\
DeepSeek V4 Flash & 50/56 (89.3\%) & 51/76 (67.1\%) & 7/18 (38.9\%) \\
gpt-5.6-luna (OpenLux) & 45/56 (80.4\%) & 50/76 (65.8\%) & 7/18 (38.9\%) \\
Qwen3.6-35B & 31/56 (55.4\%) & 27/76 (35.5\%) & 5/18 (27.8\%) \\
GPT-OSS 120B & 16/56 (28.6\%) & 16/76 (21.1\%) & 4/18 (22.2\%) \\
\bottomrule
\end{tabular}

\end{table}

\paragraph{Finding 4.}
Feature-lifting difficulty on Python-150 varies with the in-suite hard3
construction split and, more weakly, with lift type. Entanglement level is
uniformly high and cannot be tested here.

\subsection{Extraction Quality}
\label{sec:compactness}

Compactness is reported only on functional passes
(Table~\ref{tab:compactness}). Survivor sets differ by
model, so medians are not a second ranking of the same 150 tasks. Cross-model
comparisons use paired intersections.

\begin{table*}[t]
  \caption{Pass-conditioned compactness on freeze~v2 Python-150.
  RRES is submission LOC / reference LOC. Copy-heavy and compact are
  exclusive classes on the passing subset.}
  \label{tab:compactness}
  \centering
  \begin{tabular}{lrrrrr}
\toprule
\textbf{Model} & \textbf{$n$} & \textbf{Median RRES}
  & \textbf{Median copy} & \textbf{Copy-heavy} & \textbf{Compact} \\
\midrule
DeepSeek V4 Pro & 115 & 0.993 & 0.957 & 108 & 4 \\
DeepSeek V4 Flash & 108 & 0.998 & 0.968 & 105 & 1 \\
gpt-5.6-luna (OpenLux) & 102 & 0.815 & 0.164 & 54 & 37 \\
Qwen3.6-35B & 63 & 0.975 & 0.757 & 43 & 9 \\
GPT-OSS 120B & 36 & 0.983 & 0.180 & 18 & 12 \\
\bottomrule
\end{tabular}

\end{table*}

Pro and Flash passing packages are almost entirely copy-heavy (median copy
fraction 0.96--0.97). Luna's passing set is smaller in copy fraction (median
0.16) at median RRES 0.82. On the 97 tasks both Pro and Luna pass, median
copy is 0.97 versus 0.19. On the 24 tasks all five models pass, Pro/Flash
copy remains $\approx 0.96$--$0.97$ while Luna is 0.46. Correctness and
compactness therefore separate: a Functional Pass does not imply a compact
extraction relative to the frozen reference. Official Hard-50 has no
evaluator-side \texttt{reference\_solution/}, so we do not report RRES there.

\paragraph{Finding 5.}
Correctness and compactness are distinct: a functional pass does not imply
a compact extraction relative to the frozen reference.

\subsection{Case Studies}
\label{sec:cases}

Four cases instantiate the preceding findings. We describe observable
behavior, not Hidden test identifiers.

\paragraph{Process non-delivery.}
Qwen on \texttt{babel\_\_plural\_core\_\_001} ends without a package (tool
validation). The main-table failure is real; it is not a Hidden mismatch.

\paragraph{Public drift after inspecting the repository.}
On \texttt{alembic\_\_revision\_map\_core\_\_hard3\_001}, Pro and Flash
implement \texttt{get\_revision} but treat the argument \texttt{'base'} as
the symbolic base of the map, shadowing a revision whose identifier is
literally \texttt{base}. The API is present; the contract case is not.

\paragraph{Hidden completion with an open branch.}
On \texttt{aiohttp\_\_url\_params\_core\_\_hard3\_001}, exception types and
token checks exist, yet an invalid-name path required by the public contract
still does not raise. Localization is not the residue.

\paragraph{Isolation versus copy-heavy success.}
Flash on \texttt{typer\_\_command\_parser\_core\_\_001} fails Isolation via
\texttt{forbidden\_imports}. On \texttt{blinker\_\_signal\_registry\_core\_\_001},
Pro, Flash, and Luna all pass with copy-heavy packages (RRES 9.0 / 18.4 /
4.6). Passing is not compact.

\subsection{Information Boundary and Cost (Non-Main)}
\label{sec:nonmain}

Two older measurements diagnose mechanism and are not freeze~v2 Main scores.

\paragraph{Public feedback.}
A paired ablation on 12 DeepSeek V4 Flash tasks (historical slice) mounts
\texttt{public\_tests/} while keeping Hidden private. Main is 0/12;
Public-feedback is 4/12 (Table~\ref{tab:public-feedback}). All six selected
Public failures flip Public to 1; Hidden usually does not. Executable Public
tests are a real bottleneck and not a general Hidden oracle. The slice is not
the Python-150 leaderboard.

\begin{table}[t]
  \caption{Paired Public-feedback ablation on 12 historical Flash tasks.
  Not freeze~v2 Main.}
  \label{tab:public-feedback}
  \centering
  \begin{tabular}{lr}
    \toprule
    \textbf{Condition} & \textbf{Functional Pass} \\
    \midrule
    Full-Repository / No-Hint Main & 0/12 \\
    Public-feedback & 4/12 \\
    \bottomrule
  \end{tabular}
\end{table}

\paragraph{Earliest sufficiency $T^*$.}
On replayable historical Flash passing trajectories (superseded 150+E50
suite), 138 of 145 passes have gold. Median $T^*/T_{\mathrm{total}}$ is 0.40,
with a median 0.75M tokens after a sufficient package first exists
(Table~\ref{tab:tstar}). Post-sufficiency spend is mostly self-testing the
agent cannot ground in Hidden. This is not a stopping rule and is not
recomputed on freeze~v2.

\begin{table}[t]
  \caption{Earliest functional sufficiency on historical DeepSeek V4 Flash
  passing trajectories (not freeze~v2).}
  \label{tab:tstar}
  \centering
  \begin{tabular}{lrrr}
    \toprule
    \textbf{Slice} & \textbf{$n$} & \textbf{Median $T^*/T$}
      & \textbf{Post-$T^*$ tokens} \\
    \midrule
    All gold passes & 138 & 0.40 & 0.75M \\
    Direct & 58 & 0.36 & 0.74M \\
    Adapted & 52 & 0.40 & 0.79M \\
    Composite & 28 & 0.51 & 0.70M \\
    \bottomrule
  \end{tabular}
\end{table}

\paragraph{Process scaffolds.}
Development interventions (self-generated probes, test-first extraction,
repair loops, checkpoint replay, 2M caps) did not produce a stable Functional
gain over legal Main on the subsets where they were tried
(Table~\ref{tab:negative}). They are negative diagnostics, not methods in
the leaderboard.

\begin{table*}[t]
  \caption{Development interventions and their diagnostic outcomes.
  Not freeze~v2 Main.}
  \label{tab:negative}
  \centering
  \begin{tabularx}{\textwidth}{@{}p{0.25\textwidth}p{0.32\textwidth}X@{}}
    \toprule
    \textbf{Hypothesis} & \textbf{Interventions} & \textbf{Observed outcome} \\
    \midrule
    Better self-verification is sufficient
      & Self-generated probes, test-first extraction, verification ledger
      & No stable Functional gain; some variants cost more or regress \\
    Executable contract checklists close behavior
      & Self-Contract, Exec-Contract, spec-adversarial checks
      & Can repair Public; evaluated Hidden subsets do not improve reliably \\
    Correct intermediate packages are overwritten
      & Repair loops, Rescue+, best-so-far replay
      & No stable gain; 0/51 failed trajectories contain an earlier pass \\
    Tighter budget/context control improves quality
      & 2M cap, adaptive stop, compressed context, structure guidance
      & Token savings trade off against Functional Pass \\
    \bottomrule
  \end{tabularx}
\end{table*}

\section{Discussion}
\label{sec:discussion}

\subsection{Feature Lifting Is a Distinct Agent Capability}

Feature lifting asks a different question from issue repair. A repair agent
succeeds inside the original repository and can rely on its package layout,
resources, and test harness. A feature-lifting agent must infer which parts of
that environment are semantically necessary and reconstruct them behind a new
package boundary. It must satisfy an output contract after the original
repository is removed. This combination makes dependency closure and
modularization part of correctness rather than implementation style.

The task is also different from code localization. Finding a relevant symbol
can be necessary, but the \texttt{requests\_cache} and configuration examples
show why it is insufficient. Conversely, a behaviorally equivalent rewrite can
pass even if it copies little upstream code. \flb therefore evaluates delivered
behavior and artifact independence, not whether the agent followed one
prescribed extraction strategy.

\subsection{Correctness and Compactness Should Remain Separate}

A broad vendoring solution may preserve behavior but offer little reusable
modularization. A very small package may omit rare behavior. Combining these
dimensions into one scalar would obscure both failure modes and make tradeoffs
difficult to interpret. \flb treats Functional Pass as the eligibility gate
for compactness analysis. \rres is then a descriptive reference-relative
measure, not a proof of minimality or semantic authenticity.

The earlier External-50 expansion demonstrates the importance of this
separation. High pass rates alone suggested that the task was nearly solved,
while pass-conditioned footprints revealed that many solutions were close to
whole-repository copying. Hard-50 and its copy-trap tasks retain functional
difficulty while making broad copying visible.

\subsection{Implications for Agent Design}

The negative method results do not imply that agent workflows are irrelevant.
They suggest that a useful workflow must improve the evidence available for
contract closure rather than merely add procedural steps. Promising directions
include explicit obligation tracking tied to public requirements, generation
of boundary cases from exception and state semantics, provenance-aware
dependency closure, artifact-level completeness checks, and calibrated
uncertainty when a requirement cannot be verified legally.

The Public-feedback ablation suggests a layered view of feedback. Executable
Public tests can expose basic mismatches efficiently. Deeper Hidden behavior
remains a generalization problem under the Main information boundary. A robust
agent must use repository evidence and the public specification to anticipate
unobserved but declared combinations without hard-coding benchmark examples.

\subsection{Scope of the Claims}

Our main scores characterize model backends embedded in a fixed OpenHands
\mainprotocol configuration. They are not measurements of base models
independent of runtime, prompt, context policy, or provider behavior. If no
second-runtime experiment is completed, all model and trajectory claims are
explicitly scoped to this protocol. Likewise, task-factor and semantic-failure
analyses are descriptive and do not establish causal effects of repository
properties.

Harness-Bench argues that executable-agent capability should be reported at the
model--harness configuration level~\cite{yao2026harnessbench}. \flb adopts that
reporting discipline by fixing and attesting its Official Main runtime. Our
scientific axis is nevertheless different: we introduce a task construct and
study capability and failure within it, rather than treating harness variation
as the main independent variable.


\section{Threats to Validity and Responsible Release}
\label{sec:threats}

\paragraph{Construct validity.}
Functional Pass is a deterministic but finite operationalization of behavioral
completeness. Public and Hidden tests are mapped to one published contract, but
finite tests cannot prove semantic equivalence. The freeze~v2 200/0 labels
measure whether Hidden uses undeclared members and whether declared entry
points resolve; they are not a semantic Hidden-fairness gold set. Functional
Pass is binary and does not give headline partial credit. \rres, copied-code
fraction, files, and dependency footprint are proxies that may miss semantic
copying or penalize legitimate alternative implementations.

\paragraph{Internal validity.}
Results depend on source materialization, locked dependencies, prompt and
context policy, agent/evaluator containers, model endpoints, and runner state.
The reported Python-150 matrix uses freeze~v2 and the attested
\texttt{python200-prime-212930ea} images. We do not mix predecessor-freeze
agent scores into this table. Method pilots differ in subset and date; they
support mechanism hypotheses but are not pooled into a causal estimate.

\paragraph{External validity.}
Python-200$'$ is a maintainer- and AI-assisted selection of Python libraries,
tooling, and framework components. It is not a random sample of software
repositories and does not establish performance on large product applications,
GUI software, GPU systems, distributed services, or non-Python ecosystems.
Popular upstream projects may also appear in model training data. Repository
diversity does not guarantee balanced behavior or entanglement mechanisms.
Entanglement level is uniformly high on Python-150, so we do not claim an
entanglement effect.

\paragraph{Reliability and statistical conclusion validity.}
The main table uses one attempt per task and does not estimate stochastic
run-to-run variance. We report binomial confidence intervals for aggregate
rates and exact McNemar tests for paired comparisons. Neither replaces
repeated-run sensitivity. Model APIs and serving stacks may change despite
stable model names. We record freeze identity, container digests, and run
dates.

\paragraph{Annotation validity.}
Task taxonomy and semantic failure labels include AI-assisted curation.
Section~\ref{sec:mechanism} is an assistant L1 close-read, not independent
human dual review, and is not reported as gold inter-rater agreement.
Unknown rates are disclosed when they occur.

\paragraph{Runtime dependence.}
Official Main fixes OpenHands. Scores may change under another runtime, tool
interface, context policy, or recovery strategy. DeepSeek Harness and Codex
adapters exist but are not paired into the model leaderboard.

\paragraph{Licensing and release.}
Upstream repositories retain heterogeneous licenses. Release packaging, source
acquisition, evaluator-side references, and generated submissions must preserve
license obligations. Large source and result archives may be distributed via
checksummed acquisition instructions rather than committed directly.

\section{Conclusion}
\label{sec:conclusion}

We presented \flb, a benchmark for repository-level, behavior-preserving
feature extraction by code agents. The benchmark asks agents to turn a complete
but entangled repository into an independent package under a public behavioral
contract, without source-location hints or access to benchmark tests. Its
deterministic evaluator separates functional correctness from the compactness
of passing solutions.

On freeze~v2 Python-150 under Official Main, five backends span 24.0\% to
76.7\% Functional Pass. The suite is not saturated. Failures concentrate at
Public and Hidden behavior. On the two strongest models, artifact-level
failures are dominated by incomplete contract recovery rather than by missing
packages or uninspected source trees. Passing is not the same as compact
extraction. By releasing tasks, source provenance, evaluators, traces, and
reproducibility metadata, \flb aims to support research on agents that extract
reusable software rather than only modify it in place.


\bibliographystyle{ACM-Reference-Format}
\bibliography{references}

\appendix
\section{Declaration of LLM Usage}
\label{app:llm}

Large language models assisted with task curation, annotation, analysis code,
statistical summarization, and manuscript editing. All reported scientific
claims, task contracts, evaluator behavior, quantitative summaries, and final
conclusions are reviewed by the authors. AI-assisted semantic labels are marked
as such and are not reported as independent human gold.

\section{Models, Runtime, and Reproducibility Configuration}
\label{app:configuration}

Headline evaluation: freeze~v2 Python-150, OpenHands Official Main, one
attempt per task, 2026-09-04/05. Functional Pass $=$ build $\land$ public
$\land$ hidden $\land$ isolation. Empty submissions count as failures.

\begin{itemize}
  \item Freeze ID:
        \texttt{6c20ff0307762503a73cbb9ff32e9992c6446e4b17483a68373027be58cbf419}
  \item Candidate ID:
        \texttt{212930ea5363f21824afd5454c4da125052ad7a7d7186886e3dddef145811254}
  \item Agent image:
        \texttt{featureliftbench-agent:python200-prime-212930ea}\\
        \texttt{sha256:0a05b2e797e1e39ee0630131cfd9e670403ce684cd74c8375ecf7e8d444f7595}
  \item Evaluator image:
        \texttt{featureliftbench-eval:python200-prime-212930ea}\\
        \texttt{sha256:bacb3078db5e9558e9953575b5eb52f247b112dc63b7705ad7c48a7b4d864cae}
  \item Envelope: 120 steps, 128k context, No-Hint, Public/Hidden tests withheld
  \item Timeout: 3600s per task
  \item Pro suite:
        \texttt{experiments/python/openhands/deepseek-v4-pro/python150-prime-v2-main-r1}
  \item Other backends:
        \texttt{experiments/python/openhands/<model>/python200-prime-v2-main-r1},
        restricted to the 150 Python-150 task IDs
\end{itemize}

Paired McNemar tests for the headline matrix are in
Table~\ref{tab:python150-v2-pairwise}.

\begin{table}[t]
  \caption{Exact McNemar tests on freeze~v2 Python-150 (discordant tasks).}
  \label{tab:python150-v2-pairwise}
  \centering
  \begin{tabular}{llrrr}
\toprule
\textbf{A} & \textbf{B} & \textbf{A wins} & \textbf{B wins} & \textbf{McNemar $p$} \\
\midrule
Pro & Flash & 10 & 3 & 0.092 \\
Pro & Luna & 18 & 5 & 0.011 \\
Pro & Qwen & 52 & 0 & $<10^{-15}$ \\
Pro & OSS & 82 & 3 & $<10^{-20}$ \\
Flash & Luna & 16 & 10 & 0.33 \\
Flash & Qwen & 47 & 2 & $<10^{-11}$ \\
Qwen & OSS & 38 & 11 & $1.4\times10^{-4}$ \\
\bottomrule
\end{tabular}

\end{table}

\section{Official Hard-50}
\label{app:hard50}

Table~\ref{tab:hard50} reports Functional Pass on the official Hard-50
expansion and the corresponding 200-task totals. Pro was not run on Hard-50.
Hard-50 packages have no evaluator-side \texttt{reference\_solution/}, so we
do not report RRES. These rates are not the paper's difficulty claim.

\begin{table}[t]
  \caption{Official Hard-50 Functional Pass on freeze~v2 Official Main,
  with Python-150 and 200-task totals. Pro is 150-only.}
  \label{tab:hard50}
  \centering
  \begin{tabular}{lrrr}
\toprule
\textbf{Model} & \textbf{Hard-50} & \textbf{Python-150} & \textbf{Python-200$'$} \\
\midrule
DeepSeek V4 Flash & 49/50 & 108/150 & 157/200 \\
gpt-5.6-luna (OpenLux) & 42/50 & 102/150 & 144/200 \\
Qwen3.6-35B & 23/50 & 63/150 & 86/200 \\
GPT-OSS 120B & 25/50 & 36/150 & 61/200 \\
\bottomrule
\end{tabular}

\end{table}

\section{Contract-Completeness of Freeze~v2}
\label{app:violators}

The published freeze~v2 labels are 200/0/0
(Table~\ref{tab:standard-labels}, Section~\ref{sec:standard-subset}).
Table~\ref{tab:violators} lists the 32 predecessor-freeze surface or
entry-point defects that were repaired before freeze~v2; they are historical
findings, not the published labels. C4 overlap was advisory; six overlapping
Hidden cases were rewritten before freeze~v2
(Table~\ref{tab:c4-advisory}).

\begin{longtable}{llp{0.42\textwidth}}
\caption{Predecessor freeze \texttt{474862c2}: 32 tasks with a confirmed contract-completeness violation. Repaired before freeze~v2; not the published 200/0 labels.} \\
\label{tab:violators} \\
\toprule
\textbf{Task} & \textbf{Rules} & \textbf{Undeclared member or dangling symbol} \\
\midrule
\endfirsthead
\toprule
\textbf{Task} & \textbf{Rules} & \textbf{Undeclared member or dangling symbol} \\
\midrule
\endhead
aiohttp\_\_url\_params\_core\_\_hard3\_001 & C1, C2 & CIMultiDict.\_\_getitem\_\_; CIMultiDict.\_\_setitem\_\_; aiohttp.helpers.build\_url \\
apispec\_\_plugin\_documenter\_core\_\_001 & C1 & APISpec.components \\
asttokens\_\_token\_annotate\_core\_\_001 & C1 & ASTTokens.tree \\
authlib\_\_oauth2\_server\_core\_\_001 & C1 & OAuth2Request.payload \\
beaker\_\_session\_cache\_core\_\_001 & C1 & Session.\_\_setitem\_\_; Session.get; Session.id \\
build\_\_pyproject\_backend\_core\_\_hard3\_001 & C2 & build.\_builder.parse\_build\_system\_table \\
cachetools\_\_cache\_eviction\_core\_\_001 & C1 & LFUCache.\_\_contains\_\_; LFUCache.\_\_getitem\_\_; LFUCache.\_\_setitem\_\_; LRUCache.\_\_contain... \\
click\_\_lazy\_command\_core\_\_hard3\_001 & C2 & click.core.LazyCommandCollection \\
configobj\_\_roundtrip\_config\_core\_\_001 & C1 & ConfigObj.\_\_getitem\_\_ \\
cookiecutter\_\_repo\_finder\_core\_\_hard3\_001 & C2 & cookiecutter.repository.RepoFinder \\
dateutil\_\_zone\_resolver\_core\_\_hard3\_001 & C2 & dateutil.zoneinfo.ZoneResolver \\
deepdiff\_\_deep\_compare\_core\_\_001 & C1 & DeepDiff.\_\_contains\_\_ \\
diskcache\_\_eviction\_policy\_core\_\_hard3\_001 & C2 & diskcache.core.EvictionPolicyPlanner \\
fs\_\_url\_opener\_core\_\_hard3\_001 & C2 & fs.opener.registry.FSOpenerRegistry \\
hatch\_\_project\_metadata\_core\_\_hard3\_001 & C2 & hatchling.metadata.core.normalize\_project\_metadata \\
importlib\_metadata\_\_entry\_points\_core\_\_001 & C1 & EntryPoints.\_\_getitem\_\_ \\
installer\_\_wheel\_record\_core\_\_hard3\_001 & C2 & installer.records.parse\_wheel\_record \\
intervaltree\_\_interval\_tree\_core\_\_001 & C1 & IntervalTree.\_\_contains\_\_; IntervalTree.\_\_getitem\_\_ \\
jsonpointer\_\_resolve\_core\_\_001 & C1 & JsonPointer.\_\_contains\_\_ \\
multidict\_\_multidict\_mutation\_core\_\_hard3\_001 & C1 & CIMultiDict.\_\_getitem\_\_; MultiDict.\_\_getitem\_\_ \\
oslo\_config\_\_opt\_group\_core\_\_001 & C1 & ConfigOpts.host; ConfigOpts.port; ConfigOpts.timeout; ConfigOpts.worker \\
packaging\_\_requirement\_marker\_specifier\_\_001 & C1 & SpecifierSet.\_\_contains\_\_ \\
python\_configuration\_\_layered\_config\_core\_\_001 & C1 & Configuration.\_\_getitem\_\_; Configuration.y; ConfigurationSet.left; ConfigurationSet.r... \\
python\_frontmatter\_\_roundtrip\_core\_\_001 & C1 & Post.\_\_getitem\_\_ \\
readme\_renderer\_\_content\_type\_core\_\_hard3\_001 & C2 & readme\_renderer.render\_readme \\
setuptools\_scm\_\_version\_normalize\_core\_\_hard3\_001 & C2 & setuptools\_scm.version.version\_from\_scm \\
sortedcontainers\_\_sorted\_list\_core\_\_001 & C1 & SortedList.\_\_delitem\_\_ \\
spiffworkflow\_\_bpmn\_engine\_core\_\_001 & C1 & BpmnWorkflow.data \\
stevedore\_\_extension\_manager\_core\_\_hard3\_001 & C1 & ExtensionManager.\_\_getitem\_\_ \\
virtualenv\_\_interpreter\_spec\_core\_\_hard3\_001 & C2 & virtualenv.discovery.py\_info.parse\_spec \\
webob\_\_request\_response\_core\_\_001 & C1 & Response.body; Response.content\_type; Response.json\_body; Response.status\_code \\
websockets\_\_handshake\_parse\_core\_\_001 & C1 & Headers.\_\_getitem\_\_ \\
\bottomrule
\end{longtable}


\begin{table}[t]
  \caption{C4 test-overlap: predecessor advisory counts versus freeze~v2
  after the six Hidden rewrites.}
  \label{tab:c4-advisory}
  \centering
  \begin{tabular}{lrr}
\toprule
\textbf{C4 / test-overlap} & \textbf{$n$} & \textbf{Share} \\
\midrule
Predecessor freeze advisory overlaps & 29 & 14.5\% \\
Hidden cases rewritten before freeze~v2 & 6 & -- \\
Freeze~v2 static C4 hits & 0 & 0.0\% \\
\bottomrule
\end{tabular}

\end{table}

\section{Negative Intervention Registry}
\label{app:negative-registry}

The interventions in Table~\ref{tab:negative} differ in subset and date.
They are diagnostic and are not pooled into an effect size or mixed with
the freeze~v2 leaderboard.


\end{document}
